%---------------------------------------------------------------------------------------

\chapter{Implementação Computacional}\label{ch:implementacao-computacional}

%---------------------------------------------------------------------------------------

% Lista de abreviaturas e de siglas
\criarsigla{SIMD}{Single instruction, multiple data}
\criarsigla{SIMT}{Single instruction, multiple threads}
\criarsigla{TDD}{Test-driven development}
\criarsigla{SoA}{Structure of Arrays}
\criarsigla{AoS}{Array of Structures}
\criarsigla{API}{Application Programming Interface}
\criarsigla{HPC}{High Performance Computing}
\criarsigla{MPI}{Message Passing Interface}
\criarsigla{CUDA}{Compute Unified Device Architecture}
\criarsigla{JIT}{Just-In-Time}
\criarsigla{VRAM}{Video Random Access Memory}

%---------------------------------------------------------------------------------------

Neste capítulo são detalhadas a arquitetura e as tecnologias empregadas no desenvolvimento do software para simulações de dinâmica de fluidos computacional (CFD) utilizando o método de Lattice Boltzmann (LBM) para escoamentos compressíveis, com foca em decisões que impactam o desempenho em arquiteturas heterogêneas, a portabilidade entre diferentes fornecedores de hardware e a facilidade de manutenção e extensão do código.

O software adota uma arquitetura orientada a objetos em C++ moderno, estruturada para abstrair as complexidades de diferentes modelos de programação de computação de alto desempenho (HPC). O núcleo da simulação é definido por uma interface comum que estabelece o contrato para os diversos modelos de paralelização, permitindo que um orquestrador central instancie a estratégia adequada em tempo de execução com base na configuração do usuário.

A arquitetura central baseia-se no padrão de projeto \textit{Strategy}, em que uma classe controladora gerencia o ciclo de vida da simulação, a entrada e saída de dados, e a medição de desempenho.
Esta abstração permite que a lógica de controle, como a marcha no tempo apresentada no Algoritmo~\ref{alg:marcha-tempo-detalhado} e a reconstrução de campos macroscópicos para exportação, permaneça independente do hardware utilizado, seja ele uma CPU sequencial, uma GPU ou um \textit{cluster} distribuído.
\begin{algorithm}[ht]
    \caption{Marcha no Tempo para LBM Compressível com Acoplamento Térmico}\label{alg:marcha-tempo-detalhado}
    \Entrada{Dados do fluido, malha computacional, tempo final $t_{end}$}
    \Saida{Campos macroscópicos $\rho, \mathbf{u}, T$ em cada passo}
    Inicializar campos $\rho, \mathbf{u}, T$ e distribuições $f_i, g_i$\;
    \Enquanto{$t < t_{end}$}{
        \For{cada célula da malha}{
            Reconstruir $(\rho, \mathbf{u}, T)$ a partir de $f_i$ e $g_i$\;
            \tcp{Newton-Raphson para o modelo térmico}
            Resolver iterativamente para as variáveis macroscópicas\;
            Calcular equilíbrios $f_i^{eq}(\rho, \mathbf{u}, T)$ e $g_i^{eq}(\rho, \mathbf{u}, T)$\;
            Aplicar colisão (BGK): $f_i^* = f_i - \omega_f (f_i - f_i^{eq})$ e $g_i^* = g_i - \omega_g (g_i - g_i^{eq})$\;
        }
        \Se{estratégia de paralelismo é MPI}{
            Sincronizar \textit{halos} de $f_i^*$ e $g_i^*$ entre processos\;
        }
        \For{cada célula da malha}{
            Propagar (\textit{streaming}): $f_i(\mathbf{x} + \mathbf{c}_i \Delta t, t + \Delta t) = f_i^*(\mathbf{x}, t)$\;
            $g_i(\mathbf{x} + \mathbf{c}_i \Delta t, t + \Delta t) = g_i^*(\mathbf{x}, t)$\;
            Aplicar condições de contorno (paredes, entrada/saída)\;
        }
        $t \leftarrow t + \Delta t$\;
        \Se{passo de gravação}{
            Exportar resultados $(\rho, \mathbf{u}, T)$ em formato VTK\;
        }
    }
\end{algorithm}

O processo de compilação detecta automaticamente a presença de compiladores e bibliotecas como o NVIDIA HPC SDK (para OpenACC e OpenMP), NVCC (para CUDA), e implementações de MPI e OpenCL. Isso garante que o simulador possa ser compilado e executado em uma ampla variedade de ambientes.
O controle e configuração da simulação é feito através de um arquivo de configuração, onde se definem as propriedades do fluido, a malha computacional e a tecnologia de execução.

%---------------------------------------------------------------------------------------

\section{Arquitetura do Código}\label{sec:arquitetura-simulador}

A arquitetura do código foi projetada para suportar a execução em múltiplas tecnologias de paralelização sem a necessidade de reescrever a lógica física do modelo LBM. Para alcançar este objetivo, o sistema utiliza o padrão de projeto \textit{Strategy}, centralizado na interface abstrata \texttt{ILbmTube} e na classe orquestradora \texttt{Executor}.

A interface \texttt{ILbmTube} define o contrato que todos os resolvedores devem implementar, incluindo métodos para inicialização (\texttt{Init}), execução de passos de tempo (\texttt{Step} e \texttt{Run}) e acesso aos campos macroscópicos (\texttt{P}, \texttt{Rho}, \texttt{T}, \texttt{U}). Esta abstração permite que o \texttt{Executor} gerencie o ciclo de vida da simulação de forma agnóstica ao \textit{hardware}, conforme apresentado no Código~\ref{lst:interface-ilbmtube}.
\begin{lstlisting}[language=C++, caption={Interface base ILbmTube que define o contrato para os resolvedores.}, label={lst:interface-ilbmtube}, float=ht]
template<typename Scalar, typename LatticeType>
class ILbmTube {
public:
  virtual ~ILbmTube() = default;
  virtual Tensor<Scalar, LatticeType::Dim()> P() const = 0;
  virtual Tensor<Scalar, LatticeType::Dim()> Rho() const = 0;
  virtual Tensor<Scalar, LatticeType::Dim()> T() const = 0;
  virtual Tensor<Scalar, LatticeType::Dim() + 1> U() const = 0;
  virtual void Init() = 0;
  virtual void Step(bool save) = 0;
  virtual void Run(Index steps, bool save) = 0;
};
\end{lstlisting}

O \texttt{Executor} atua como um \textit{Singleton} que seleciona a implementação correta de \texttt{ILbmTube} em tempo de execução, baseando-se no arquivo de configuração (\texttt{config.yaml}). Ele também é responsável pela medição de desempenho e exportação de resultados em formato VTK para visualização posterior.
O impacto desta arquitetura na complexidade do código é positivo, pois isola a complexidade das APIs de paralelismo em classes específicas, facilitando a inclusão de novas estratégias de paralelização.

Um componente crítico para a performance é a classe \texttt{DeviceBuffer}, que gerencia a alocação de memória.
Em arquiteturas de CPU, ela se comporta como um \texttt{std::vector} alinhado, enquanto em GPUs ela encapsula ponteiros de memória de dispositivo (\textit{device pointers}), residindo no que se denomina memória do dispositivo (\textit{device memory}). A escolha do \textit{layout} de memória Structure of Arrays (SoA) é fundamental: ao armazenar cada componente da distribuição de probabilidade ($f_0, f_1, \dots, f_8$) em blocos contíguos de memória, maximiza-se o uso do cache em CPUs através da busca antecipada (\textit{prefetching}), mecanismo no qual o processador carrega dados da memória principal para o cache antes que sejam explicitamente solicitados.
Além disso, este \textit{layout} garante o coalescimento de memória (\textit{memory coalescing}) em GPUs, técnica que permite combinar múltiplos acessos à memória global realizados por diferentes \textit{threads} em uma única transação, desde que os endereços acessados sejam adjacentes simultaneamente.

A marcha no tempo segue uma estrutura de pipeline que minimiza transferências entre a memória do sistema principal (\textit{host memory}), referente à RAM do sistema, e a memória do dispositivo (\textit{device memory}), referente à VRAM da GPU, conforme orquestrado pelo ciclo principal.

\section{Implementação Computacional Sequencial}\label{sec:implementacao-computacional-sequencial}

A versão sequencial, referida no código como \textit{Plain}, estabelece a referência funcional para o projeto.
Ela implementa o modelo LBM compressível seguindo a formulação de \textcite{Tran2022}, utilizando as classes de tensores da biblioteca Eigen para manipulação de dados multidimensionais em C++.

O objetivo desta implementação, além de fornecer uma referência para validação científica, é servir como um modelo de arquitetura para as demais implementações.
Por esta razão, o código evita construções excessivamente específicas de CPU, como o uso de listas ligadas ou ponteiros complexos, preferindo uma estrutura de dados linear e laços de repetição simples que podem ser facilmente portados para outras tecnologias HPC\@.

A estrutura principal de dados utiliza o \textit{layout} SoA (\textit{Structure of Arrays}), onde as distribuições de probabilidade $f_i$ e $g_i$ para cada uma das $Q=9$ velocidades da rede D2Q9 são armazenadas em buffers contíguos.
Esta escolha visa facilitar a vetorização por parte do compilador e serve como base para o acesso coalescido de memória nas variantes de aceleradores.
O algoritmo principal é dividido em etapas discretas encapsuladas na classe \texttt{lbmini::plain::LbmTube}:

\begin{enumerate}
    \item \textbf{Cálculo de grandezas macroscópicas}: reconstrução de $\rho$, $\mathbf{u}$, $T$ e $p$ a partir dos momentos de $f_i$ e $g_i$, em que $\mathbf{u}$ é o vetor velocidade e $p$ a pressão local.
    \item \textbf{Colisão e equilíbrio}: aplicação do operador BGK e cálculo das distribuições de equilíbrio $f_i^{eq}$ e $g_i^{eq}$.
    Esta etapa é a mais custosa computacionalmente devido à solução iterativa utilizando Newton-Raphson.
    \item \textbf{Streaming}: transporte das distribuições entre nós vizinhos da malha utilizando interpolação por pesos de distância inversa (\textit{Inverse Distance Weighting}, IDW).
\end{enumerate}

A etapa de colisão é dividida em dois subproblemas: a determinação do equilíbrio para o campo de densidade ($f_i^{eq}$) e para o campo de energia total ($g_i^{eq}$). Enquanto $f_i^{eq}$ é calculado via forma de produto fatorado, a determinação de $g_i^{eq}$ exige a solução de um sistema não linear para garantir a conservação de energia em regimes compressíveis.
Na implementação computacional, isto é feito através de um laço fixo de iterações Newton-Raphson, conforme mostrado no Código~\ref{lst:plain-newton}.
Esta escolha de laço fixo é estratégica para garantir a previsibilidade do fluxo de execução em futuras implementações SIMT para GPUs.
\begin{lstlisting}[language=C++, caption={Resolvedor Newton-Raphson para equilíbrio térmico.}, label={lst:plain-newton}, float=ht]
for (Index iter = 0; iter < maxIter; ++iter) {
    // Calculo da Jacobiana J e do residuo r
    const Scalar detJ = J[0][0] * J[1][1] - J[0][1] * J[1][0];
    const Scalar invDet = Scalar{ 1 } / detJ;

    // Atualizacao dos multiplicadores de Lagrange xi
    xi[0] -= alpha * (J[1][1] * r[0] - J[0][1] * r[1]) * invDet;
    xi[1] -= alpha * (-J[1][0] * r[0] + J[0][0] * r[1]) * invDet;
    alpha *= Scalar{ 0.5 }; // Reducao do passo
}
\end{lstlisting}
Em que $\xi$ representa os multiplicadores de Lagrange, $\alpha$ é o fator de amortecimento do passo, $J$ denota a matriz Jacobiana e $r$ indica o vetor resíduo.

No simulador, a implementação sequencial utiliza laços aninhados que percorrem a malha, onde cada célula é processada de forma independente.
O Código~\ref{lst:plain-collision} exemplifica a estrutura de processamento da colisão na CPU, evidenciando o uso do \textit{layout} SoA para acesso aos dados.
\begin{lstlisting}[language=C++, caption={Processamento sequencial da colisão e equilíbrio.}, label={lst:plain-collision}, float=ht]
for (Index cell = 0; cell < N; ++cell) {
    const Scalar rho = pRho[cell];
    const Scalar Tv = pT[cell];
    const Scalar ux = pU[0 * N + cell];
    const Scalar uy = pU[1 * N + cell];

    // Calculo de feq e geq (Newton-Raphson omitido para brevidade)
    compute_equilibria(cell, rho, Tv, ux, uy, feq, geq);

    // Aplicacao do operador BGK em cada direcao idc
    for (Index idc = 0; idc < 9; ++idc) {
        const Index idx = idc * N + cell;
        pF[idx] = pF[idx] + omegaL * (feq[idc] - pF[idx]);
        pG[idx] = pG[idx] + omegaL * (geq[idc] - pG[idx]);
    }
}
\end{lstlisting}
Em que $N$ é o número total de células, \texttt{cell} é o índice da célula atual, \texttt{idc} representa a direção da velocidade discreta, $\omega_L$ é o parâmetro de relaxação e $f_i$, $g_i$ são as funções de distribuição para densidade e energia, respectivamente.

A etapa de \textit{streaming} é implementada utilizando tabelas de índices pré-computadas, que mapeiam os vizinhos de cada célula considerando as condições de contorno periódicas ou de parede.
Isso elimina a necessidade de cálculos geométricos repetitivos durante a marcha no tempo.

%---------------------------------------------------------------------------------------

\section{Implementação Computacional Utilizando OpenMP}\label{sec:implementacao-computacional-utilizando-openmp}
%---------------------------------------------------------------------------------------

A implementação OpenMP suporta tanto paralelismo em multicore quanto aceleração de \textit{hardware} com GPUs.
\begin{lstlisting}[language=C++, caption={Diretiva de paralelização OpenMP para CPUs multinúcleo.}, label={lst:openmp-cpu}, float=ht]
#pragma omp parallel for schedule(static)
for (Index cell = 0; cell < N; ++cell) {
    // Implementacao da colisao e streaming
    collide_and_stream(cell);
}
\end{lstlisting}
Para CPUs, utiliza-se a diretiva mostrada no Código~\ref{lst:openmp-cpu} para garantir que os laços de células sejam distribuídos entre núcleos e vetorizados via instruções SIMD \parencite{Dagum1998}, onde $N$ representa o número total de células e \texttt{cell} o índice da iteração.

A arquitetura para OpenMP foca na escalabilidade vertical (\textit{strong scaling}) em sistemas com múltiplos processadores físicos.
Um fator crítico de performance abordado é a afinidade de \textit{threads} e o acesso à memória em sistemas NUMA (\textit{Non-Uniform Memory Access}). Ao utilizar a cláusula \texttt{schedule(static)}, o simulador garante que as mesmas células sejam processadas pelas mesmas \textit{threads} em cada passo de tempo, o que melhora a localidade de dados no cache L3 e minimiza o tráfego no barramento entre processadores.
\begin{lstlisting}[language=C++, caption={Uso de SIMD e afinidade no OpenMP.}, label={lst:openmp-simd}, float=ht]
#pragma omp parallel for simd schedule(static)
for (Index cell = 0; cell < N; ++cell) {
    // O compilador pode usar instrucoes AVX/AVX-512 aqui
    pRho[cell] = 0;
    for (int idc = 0; idc < 9; ++idc) {
        pRho[cell] += pF[idc * N + cell];
    }
}
\end{lstlisting}
Nas implementações em paralelo, a variável \texttt{cell} é utilizada como índice para a distribuição de carga de trabalho entre os núcleos ou \textit{threads} de execução.

A paralelização é configurada com escalonamento estático para minimizar a latência de gerência de \textit{threads} em uma malha uniforme.
Na versão para aceleradores de \textit{hardware}, a gestão de memória é feita de forma explícita via diretivas \texttt{\#pragma omp target data}, garantindo que os dados residam no acelerador durante toda a simulação e evitando o custo proibitivo de transferência via comunicação PCIe entre passos de tempo.

A diretiva \texttt{collapse(2)} é aplicada em malhas bidimensionais para aumentar o espaço de iteração disponível para o escalonador da GPU. Além disso, utiliza-se a cláusula \texttt{thread\_limit} para sintonizar a ocupação dos multiprocessadores.
O impacto na complexidade é moderado, pois as diretivas OpenMP são menos intrusivas que o CUDA ou OpenCL, permitindo manter o código C++ muito próximo da versão sequencial.
No entanto, a sintonização fina para performance (\textit{performance tuning}) exige um conhecimento profundo do comportamento do compilador e da arquitetura alvo.

%---------------------------------------------------------------------------------------

\section{Implementação Computacional Utilizando OpenACC}\label{sec:implementacao-computacional-utilizando-openacc}
%---------------------------------------------------------------------------------------

A tecnologia OpenACC foca na produtividade e portabilidade entre diferentes famílias de GPUs.
A lógica de implementação é muito próxima da versão sequencial e da versão utilizando OpenMP, mas decorada com diretivas \texttt{\#pragma acc} que instruem o compilador sobre o paralelismo disponível e a gestão de dados entre as memórias do sistema principal (\textit{host memory}) e do dispositivo acelerador (\textit{device memory}) \parencite{Wolfe2018}.

O impacto do OpenACC na complexidade do código é mínimo, pois permite manter uma única base de código para CPU e GPU. A implementação computacional utiliza o modelo de dados estruturados do OpenACC para garantir que grandes tensores de distribuições $f$ e $g$ sejam mantidos na memória da GPU durante todo o laço de tempo.
Utiliza-se a cláusula \texttt{present} para indicar que os dados já foram movidos para o dispositivo, evitando transferências redundantes no barramento PCIe que degradariam severamente o desempenho.

A estruturação do código em funções \textit{inline} marcadas com \texttt{\#pragma acc routine seq} permite que a lógica complexa do método Newton-Raphson seja reutilizada sem duplicação de código, mantendo a modularidade.
Esta é uma vantagem significativa em termos de manutenção em comparação com o CUDA, onde a lógica precisaria ser duplicada ou movida para arquivos de cabeçalho com qualificadores de dispositivo (\textit{device}).
A paralelização em OpenACC é expressa de forma descritiva, onde o desenvolvedor define as regiões paralelas e o compilador mapeia as iterações para os aceleradores, conforme ilustrado no Código~\ref{lst:openacc-parallel}.
Para otimizar a performance, utiliza-se a cláusula \texttt{vector\_length(256)} para sintonizar o tamanho do vetor de execução na GPU, além de garantir que os laços internos de velocidades sejam processados de forma sequencial por cada \textit{thread} (\texttt{routine seq}), maximizando a reutilização de registradores.
\begin{lstlisting}[language=C++, caption={Uso de acc routine para reutilização de código.}, label={lst:openacc-routine}, float=ht]
#pragma acc routine seq
void compute_feq_acc(Scalar rho, Scalar u, int idc) {
    // Mesma logica da versao sequencial
    return ...;
}
\end{lstlisting}
\begin{lstlisting}[language=C++, caption={Diretivas OpenACC para paralelização em GPU.}, label={lst:openacc-parallel}, float=ht]
#pragma acc parallel loop present(pF, pG, pRho, pT, pU)
for (int cell = 0; cell < N; ++cell) {
    const Scalar rho = pRho[cell];
    // Execucao paralela no dispositivo...
}
\end{lstlisting}

%---------------------------------------------------------------------------------------

\section{Implementação Computacional Utilizando CUDA}\label{sec:implementacao-computacional-utilizando-cuda}
%---------------------------------------------------------------------------------------

A implementação em CUDA representa a versão de máxima performance para arquiteturas NVIDIA. Diferente da versão sequencial, os kernels CUDA são especializados para explorar o paralelismo massivo de nível de \textit{thread} e a hierarquia de memória da GPU \parencite{Chakrabarti2012}.

A escolha de arquitetura no CUDA foca em dois pilares: a maximização da ocupação dos multiprocessadores (\textit{Streaming Multiprocessors}, SMs) e a eficiência do subsistema de memória.
Para isso, o simulador mapeia cada célula da malha para uma \textit{thread} CUDA individual.
Dada a natureza do algoritmo LBM, que é limitado por largura de banda de memória (\textit{bandwidth-bound}), o uso do layout SoA é crucial para permitir que as requisições de memória de um \textit{warp} sejam combinadas em uma única transação de memória global (\textit{coalescing}).

Uma das principais otimizações é o uso de memória constante para armazenar as velocidades discretas da rede ($c_x, c_y$) e os pesos.
Como estes dados são lidos simultaneamente por todas as \textit{threads} de um \textit{warp}, o uso de cache constante reduz a latência de acesso em comparação com a memória global.
Além disso, utiliza-se a memória compartilhada (\textit{shared memory}) em etapas específicas para reduzir o tráfego com a memória global (\textit{VRAM}), embora a implementação atual priorize o uso de registradores devido à baixa complexidade geométrica da malha.

A implementação utiliza \textit{double buffering} para as distribuições $f$ e $g$, eliminando dependências de escrita durante o \textit{streaming}.
Além disso, emprega-se a técnica de \textit{loop unrolling} manual via \texttt{\#pragma unroll} nos laços internos de velocidades ($Q=9$), o que permite ao compilador \texttt{nvcc} otimizar o uso de registradores e reduzir o custo de controle de fluxo.
A paralelização é executada distribuindo as células da malha entre os blocos e \textit{threads} da GPU.
O Código~\ref{lst:cuda-parallel} detalha como o índice global da célula é mapeado para a hierarquia de execução CUDA e como o acesso aos dados segue o padrão SoA para garantir o coalescimento de memória (\textit{memory coalescing}), em que \texttt{blockIdx}, \texttt{blockDim} e \texttt{threadIdx} são variáveis intrínsecas do CUDA para identificação de \textit{threads}.
Outro ponto relevante é o uso de fluxos CUDA (\textit{CUDA Streams}) para permitir a sobreposição de transferência de dados (\texttt{cudaMemcpyAsync}) com a execução de kernels.
\begin{lstlisting}[language=C++, caption={Uso de registradores e unroll em kernel CUDA.}, label={lst:cuda-unroll}, float=ht]
// Unrolling para forcar o uso de registradores
#pragma unroll
for (int idc = 0; idc < 9; ++idc) {
    const Scalar fi = pF[idc * N + cell];
    const Scalar feq = compute_feq(rho, u, idc);
    pFaux[idc * N + cell] = fi + omega * (feq - fi);
}
\end{lstlisting}
\begin{lstlisting}[language=C++, caption={Lógica de paralelização e acesso SoA em kernel CUDA.}, label={lst:cuda-parallel}, float=ht]
__global__ void collisionKernel(Scalar* pF, Scalar* pG, int N) {
    // Mapeamento de thread para celula da malha
    int cell = blockIdx.x * blockDim.x + threadIdx.x;

    if (cell < N) {
        // Acesso SoA: idc * N + cell garante acesso contiguo entre threads
        for (int idc = 0; idc < 9; ++idc) {
            const int idx = idc * N + cell;
            Scalar fi = pF[idx];
            // Processamento local da colisao...
            pF[idx] = fi_new;
        }
    }
}
\end{lstlisting}

Para minimizar a latência de lançamento de kernels, as etapas de colisão e \textit{streaming} são agrupadas em um único kernel (\textit{kernel fusion}), reduzindo o tráfego de memória global ao manter valores intermediários em registradores entre as operações.

%---------------------------------------------------------------------------------------

\section{Implementação Computacional Utilizando OpenCL}\label{sec:implementacao-computacional-utilizando-opencl}
%---------------------------------------------------------------------------------------

O módulo OpenCL provê portabilidade funcional para diversos aceleradores, incluindo GPUs de diferentes fabricantes (NVIDIA, AMD, Intel) e FPGAs.
A implementação utiliza a API C++ de alto nível para gerenciar a complexidade do padrão OpenCL\@, abstraindo a criação de contextos, filas de comando e gerenciamento de buffers, o que reduz a verbosidade do código principal.

Diferente do CUDA, onde a paralelização é integrada ao compilador da linguagem, o OpenCL adota um modelo de plataforma e dispositivo.
A arquitetura do simulador lida com isso através da compilação JIT (\textit{Just-In-Time}). Os kernels são escritos em linguagem OpenCL C (baseada em C99) e fornecidos como cadeias de caracteres ou arquivos externos, sendo compilados apenas durante a fase de inicialização do programa.
Esta abordagem permite que o software realize otimizações específicas para o \textit{hardware} detectado no momento da execução, como o ajuste do tamanho do grupo de trabalho (\textit{work-group size}) e o uso de extensões de precisão dupla (\texttt{cl\_khr\_fp64}).

O layout de memória SoA é rigorosamente mantido para garantir acessos coalescidos, o que é vital para a performance em GPUs AMD e Intel.
\begin{lstlisting}[language=C, caption={Lógica de indexação global em kernel OpenCL.}, label={lst:opencl-kernel}, float=ht]
__kernel void collisionKernel(__global Scalar* pF, __global Scalar* pG, int N) {
    // Obtencao do indice global da celula
    int cell = get_global_id(0);

    if (cell < N) {
        // Acesso SoA e processamento...
    }
}
\end{lstlisting}
Diferente do CUDA, onde a sintaxe é integrada ao C++, no OpenCL a paralelização é expressa no kernel utilizando funções intrínsecas para obter o índice global de execução, como demonstrado no Código~\ref{lst:opencl-kernel}.
A complexidade da implementação OpenCL reside principalmente no gerenciamento explícito da sincronização entre o sistema principal e o dispositivo de aceleração de \textit{hardware} através de eventos (\texttt{cl::Event}), que são utilizados no código para enfileirar as etapas da marcha no tempo sem bloqueios desnecessários da CPU\@.
\begin{lstlisting}[language=C++, caption={Gerenciamento de buffers e kernels via API C++ do OpenCL.}, label={lst:opencl-cpp}, float=ht]
// Criacao de buffers no dispositivo
cl::Buffer bufF(context, CL_MEM_READ_WRITE, sizeInBytes);
cl::Buffer bufRho(context, CL_MEM_READ_WRITE, rhoSizeInBytes);

// Configuracao dos argumentos e despacho do kernel
kernelMacroscopic.setArg(0, bufF);
kernelMacroscopic.setArg(1, bufRho);
queue.enqueueNDRangeKernel(kernelMacroscopic, cl::NullRange, cl::NDRange(N), cl::NullRange);
\end{lstlisting}

Um desafio no OpenCL é a latência no despacho de comandos.
Para mitigar isso, a implementação utiliza filas de execução configuradas para execução fora de ordem quando suportado, e buffers de comandos para agrupar as chamadas de kernel da marcha no tempo em uma única submissão ao dispositivo, reduzindo a carga no processador principal (\textit{host processor}).
Em termos de performance, o OpenCL apresenta resultados competitivos com o CUDA, embora o custo de compilação inicial (JIT) e a necessidade de kernels em arquivos separados aumentem a complexidade de manutenção do código.

%---------------------------------------------------------------------------------------

\section{Implementação Computacional Utilizando Kokkos}\label{sec:implementacao-computacional-utilizando-kokkos}
%---------------------------------------------------------------------------------------

A arquitetura utilizando Kokkos baseia-se no uso de \textit{Views} para o gerenciamento de memória e \textit{Parallel Dispatch} para a execução dos kernels.
Ao contrário das versões nativas que exigem a gestão manual de buffers entre o hospedeiro e o dispositivo, o Kokkos utiliza espaços de memória (\textit{Memory Spaces}) e \textit{Execution Spaces} que são resolvidos em tempo de compilação para o módulo selecionado (como CUDA para GPUs NVIDIA ou OpenMP para CPUs multicore).
O impacto na complexidade é reduzido pela unificação do código de paralelismo, permitindo que a mesma lógica de colisão e \textit{streaming} seja executada em diferentes arquiteturas apenas alterando os parâmetros de compilação.

%---------------------------------------------------------------------------------------

\section{Implementação Computacional Utilizando RAJA}\label{sec:implementacao-computacional-utilizando-raja}
%---------------------------------------------------------------------------------------

A integração com o RAJA permite a escrita de laços de repetição agnósticos à plataforma através de políticas de execução (\textit{Execution Policies}).
O diferencial do RAJA reside na sua flexibilidade para compor diferentes estratégias de paralelização utilizando \textit{templates} de C++, facilitando a experimentação com diferentes níveis de paralelismo (como a combinação de RAJA com CUDA ou OpenMP) sem alterar a estrutura fundamental do código-fonte.
Estas implementações representam a vanguarda da portabilidade de desempenho, servindo como base para futuras extensões do simulador em ambientes de computação heterogênea de próxima geração.

%---------------------------------------------------------------------------------------

%---------------------------------------------------------------------------------------

\section{Implementação Computacional Utilizando MPI}\label{sec:implementacao-computacional-utilizando-mpi}
%---------------------------------------------------------------------------------------

Para permitir simulações de larga escala que excedem a capacidade de memória e processamento de um único sistema computacional (também chamado de nó de computação), utiliza-se o padrão MPI para paralelismo em uma arquitetura com memória distribuída.
A arquitetura distribuída baseia-se na decomposição de domínio unidimensional ao longo do eixo longitudinal ($x$), o que é particularmente eficiente para o estudo de tubos de choque, onde os gradientes mais intensos ocorrem nesta direção.

A complexidade da implementação MPI reside na gestão explícita da comunicação e na consistência dos dados nas fronteiras dos subdomínios presentes em cada nó de computação.
Cada processo MPI é responsável por uma fatia da malha total e mantém uma cópia local dos dados necessários para computar a colisão e o \textit{streaming}.
A sincronização entre processos é realizada através da técnica de \textit{halos} ou \textit{ghost cells}, que são camadas extras de células nas bordas de cada subdomínio que armazenam os dados pertencentes aos processos vizinhos.

Após a etapa de colisão, cada processo envia as distribuições em suas fronteiras para os vizinhos imediatos e recebe as distribuições correspondentes utilizando comunicações não-bloqueantes (\texttt{MPI\_Isend} e \texttt{MPI\_Irecv}).
Isso permite a sobreposição de comunicação e computação (\textit{latency hiding}): enquanto as distribuições de fronteira viajam pela rede, o processo continua calculando a colisão nas células internas do domínio, onde os dados não dependem de informações externas.
O impacto na performance é significativo em clusters de alta velocidade, permitindo escalabilidade quase linear (\textit{weak scaling}) para problemas de grande porte.

A paralelização via MPI exige o empacotamento explícito dos dados de fronteira em buffers contíguos antes da transmissão, como mostrado no Código~\ref{lst:mpi-pack}, onde \texttt{ghosts} indica a largura da zona de troca (\textit{halo}), \texttt{ny} é a dimensão transversal da malha, \texttt{sendL} e \texttt{recvL} são os buffers de comunicação, e \texttt{reqs} armazena os estados das operações não-bloqueantes.
Esta etapa é necessária porque os dados de interesse podem não estar contíguos na memória devido ao layout SoA bidimensional.
Além disso, o simulador utiliza tipos de dados derivados do MPI (\texttt{MPI\_Type\_vector}) para automatizar este processo em malhas mais complexas, simplificando a codificação conforme ilustrado no Código~\ref{lst:mpi-type}.
\begin{lstlisting}[language=C++, caption={Uso de tipos derivados MPI para troca de halos.}, label={lst:mpi-type}, float=ht]
// Criacao de tipo vetor para colunas de halo
MPI_Datatype haloType;
MPI_Type_vector(ny, Q * 2, nx_local, MPI_DOUBLE, &haloType);
MPI_Type_commit(&haloType);

// Envio nao-bloqueante usando o tipo customizado
MPI_Isend(&pF[0], 1, haloType, leftRank, tag, comm, &request);
\end{lstlisting}
\begin{lstlisting}[language=C++, caption={Empacotamento de halos e comunicação não-bloqueante MPI.}, label={lst:mpi-pack}, float=ht]
// Empacotamento de dados de f e g para envio
for (Index idc = 0; idc < 9; ++idc) {
    for (Index i = 0; i < ghosts; ++i) {
        for (Index j = 0; j < ny; ++j) {
            Index cell = (i + ghosts) * ny + j;
            sendL[idx++] = pF[idc * N + cell];
            sendL[idx++] = pG[idc * N + cell];
        }
    }
}

// Comunicacao nao-bloqueante (Isend/Irecv)
MPI_Isend(sendL, count, mpiType, leftRank, 0, comm, &reqs[0]);
MPI_Irecv(recvL, count, mpiType, leftRank, 1, comm, &reqs[1]);
\end{lstlisting}

A arquitetura da implementação utilizando MPI encapsula a lógica de comunicação, garantindo que o restante do código veja o sistema distribuído de forma transparente, mantendo a consistência da arquitetura \textit{Strategy}.
